\begin{proof}[Proof of \cref{obj:lem:pair-polynomial-modulus}]
\begin{enumerate}
\item \emph{Real fixed vectors.}\label{proof:pair-polynomial-modulus:fixed-vectors}
Define the constant column vector and, for an arbitrary real function, its extrema and midpoint by
\[
\mathbf1(s):=1\quad(s\in\mathcal S),\qquad
f:\mathcal S\to\mathbb R,\qquad
f_{\min}:=\min_{s\in\mathcal S}f(s),\qquad
f_{\max}:=\max_{s\in\mathcal S}f(s),\qquad
c_f:=\frac{f_{\min}+f_{\max}}2.
\]
For probability vectors \(\omega_1,\omega_2:\mathcal S\to\mathbb R\), put
\[
\tau(\omega_1,\omega_2)
:=\frac12\sum_{s\in\mathcal S}|\omega_1(s)-\omega_2(s)|.
\]
Their equal total masses give
\[
\sum_s(\omega_1(s)-\omega_2(s))=0,
\qquad
|f(s)-c_f|\le\frac{f_{\max}-f_{\min}}2.
\]
Consequently, centering at \(c_f\) and applying the triangle inequality yields
\[
\begin{aligned}
\left|\sum_s\omega_1(s)f(s)-\sum_s\omega_2(s)f(s)\right|
&=\left|\sum_s(\omega_1(s)-\omega_2(s))(f(s)-c_f)\right|\\
&\le\sum_s|\omega_1(s)-\omega_2(s)|
                 \frac{f_{\max}-f_{\min}}2\\
&=\tau(\omega_1,\omega_2)(f_{\max}-f_{\min}).
\end{aligned}
\]
Now suppose \(Pf=f\), and choose
\[
i_{\min},i_{\max}\in\mathcal S,\qquad
f(i_{\min})=f_{\min},\qquad f(i_{\max})=f_{\max}.
\]
The rows of \(P\) are probability vectors, and
\[
\tau\bigl(P(i_{\max},\cdot),P(i_{\min},\cdot)\bigr)
\le\delta(P)\le\alpha.
\]
Applying the preceding estimate to these rows gives
\[
\begin{aligned}
0\le f_{\max}-f_{\min}
&=\left|(Pf)(i_{\max})-(Pf)(i_{\min})\right|\\
&\le
\tau\bigl(P(i_{\max},\cdot),P(i_{\min},\cdot)\bigr)
       (f_{\max}-f_{\min})\\
&\le\alpha(f_{\max}-f_{\min}).
\end{aligned}
\]
Since \(1-\alpha>0\), the extrema coincide, so \(f=f_{\min}\mathbf1\).
Conversely, row-stochasticity gives \(P\mathbf1=\mathbf1\). Thus
\[
\ker(P-\mathrm{Id})=\operatorname{span}_{\mathbb R}\{\mathbf1\}.
\]
The same argument applies to \(P'\).
% lean: CausalSmith.Stat.PomdpBinaryhiddenRate.dobrushin_fixedVector_constant

\item \emph{The exact polynomial identity.}\label{proof:pair-polynomial-modulus:identity}
Write
\[
\chi_P(z):=\det(z\,\mathrm{Id}-P),\qquad
\chi_{P'}(z):=\det(z\,\mathrm{Id}-P').
\]
Because \(\mathbf1\ne0\) and \(P\mathbf1=\mathbf1\), we have
\(\chi_P(1)=0\), and similarly \(\chi_{P'}(1)=0\).
Define the polynomial quotients
\[
q_P(z):=\frac{\chi_P(z)}{z-1}\in\mathbb R[z],
\qquad
q_{P'}(z):=\frac{\chi_{P'}(z)}{z-1}\in\mathbb R[z].
\]
Thus
\[
\chi_P(z)=(z-1)q_P(z),\qquad
\chi_{P'}(z)=(z-1)q_{P'}(z),\qquad
Q_{P,P'}(z)=q_P(z)q_{P'}(z).
\]
The characteristic polynomials have degree four, so
\[
\deg q_P=\deg q_{P'}=3,\qquad \deg Q_{P,P'}\le6.
\]

For polynomial matrix evaluation, use the convention
\[
F(z)=\sum_{k=0}^{n_F}F_kz^k\in\mathbb R[z],
\qquad n_F\in\mathbb N,\qquad
F(P):=\sum_{k=0}^{n_F}F_kP^k.
\]
Stationarity and induction give
\[
\pi P^k=\pi\quad(k\in\mathbb N),
\qquad
\pi F(P)=\sum_{k=0}^{n_F}F_k\pi P^k=F(1)\pi.
\]
Cayley--Hamilton, together with the characteristic factorization, yields
\[
(P-\mathrm{Id})q_P(P)=\chi_P(P)=0.
\]
Each column of \(q_P(P)\) therefore belongs to the fixed space established above.
Since each column is constant and \(\sum_s\pi(s)=1\), for every \(i,j\in\mathcal S\),
\[
[q_P(P)]_{ij}
=\sum_s\pi(s)[q_P(P)]_{sj}
=[\pi q_P(P)]_j
=q_P(1)\pi(j).
\]
With the rank-one matrix defined by
\[
(\mathbf1\pi)_{ij}:=\pi(j)\quad(i,j\in\mathcal S),
\]
this says
\[
q_P(P)=q_P(1)\mathbf1\pi.
\]
Multiplying by \(q_{P'}(P)\) and using polynomial stationarity gives
\[
\begin{aligned}
Q_{P,P'}(P)
&=q_P(P)q_{P'}(P)\\
&=q_P(1)\mathbf1\bigl(\pi q_{P'}(P)\bigr)\\
&=q_P(1)q_{P'}(1)\mathbf1\pi
=Q_{P,P'}(1)\mathbf1\pi.
\end{aligned}
\]
Applying the same calculation to \(P'\), with the polynomial factors in the opposite order, gives
\[
Q_{P,P'}(P')=Q_{P,P'}(1)\mathbf1\pi'.
\]
% lean: CausalSmith.Stat.PomdpBinaryhiddenRate.transientPoly_mul_aeval_projection

Since the initial vectors have total mass one, these projection identities imply
\[
\nu Q_{P,P'}(P)r=Q_{P,P'}(1)\pi r,
\qquad
\nu'Q_{P,P'}(P')r'=Q_{P,P'}(1)\pi'r'.
\]
On the other hand, expanding the polynomial of degree at most six gives
\[
\nu Q_{P,P'}(P)r=\sum_{k=0}^6Q_k\nu P^kr,
\qquad
\nu'Q_{P,P'}(P')r'=\sum_{k=0}^6Q_k\nu'P'^kr'.
\]
Subtracting proves
\[
Q_{P,P'}(1)(\pi r-\pi'r')
=\sum_{k=0}^6Q_k\bigl(\nu P^kr-\nu'P'^kr'\bigr).
\]
% lean: CausalSmith.Stat.PomdpBinaryhiddenRate.pairPolynomial_value_identity

\item \emph{Algebraic simplicity at one.}
Define the shifted operators
\[
D_P:=P-\mathrm{Id},\qquad D_{P'}:=P'-\mathrm{Id}.
\]
For a column vector \(w\in\mathbb R^{\mathcal S}\) satisfying \(D_P^2w=0\), put
\[
w_{\mathrm{shift}}:=D_Pw.
\]
Then
\[
D_Pw_{\mathrm{shift}}=0,\qquad
Pw_{\mathrm{shift}}=w_{\mathrm{shift}},
\]
so \(w_{\mathrm{shift}}\) is constant by \cref{proof:pair-polynomial-modulus:fixed-vectors}.
Stationarity and normalization now give, for every \(s\in\mathcal S\),
\[
\begin{aligned}
w_{\mathrm{shift}}(s)
&=\sum_i\pi(i)w_{\mathrm{shift}}(i)\\
&=\pi D_Pw
=(\pi P-\pi)w=0.
\end{aligned}
\]
Hence
\[
D_P^2w=0\quad\Longrightarrow\quad D_Pw=0.
\]

Induction extends this implication to
\[
D_P^nw=0\quad\Longrightarrow\quad D_Pw=0
\qquad(n\in\mathbb N,\ w\in\mathbb R^{\mathcal S}).
\]
For \(n=0\), the premise gives \(w=0\).
For the induction step,
\[
D_P^{n+1}w=0
\quad\Longrightarrow\quad
D_P^n(D_Pw)=0
\quad\Longrightarrow\quad
D_P^2w=0
\quad\Longrightarrow\quad
D_Pw=0,
\]
where the middle implication uses the induction hypothesis.

Thus the generalized eigenspace at one is
\[
\mathcal E_{P,1}
:=\bigcup_{n\in\mathbb N}\ker(D_P^n)
=\ker D_P
=\operatorname{span}_{\mathbb R}\{\mathbf1\}.
\]
Algebraic multiplicity equals the dimension of this generalized eigenspace.
Since \(\mathbf1\ne0\), it follows that
\[
m_P
:=\max\{n\in\mathbb N:(z-1)^n\text{ divides }\chi_P(z)\}
=\dim_{\mathbb R}\mathcal E_{P,1}=1.
\]
The identical argument for \(P'\) gives
\[
m_{P'}
:=\max\{n\in\mathbb N:(z-1)^n\text{ divides }\chi_{P'}(z)\}
=1.
\]
% lean: CausalSmith.Stat.PomdpBinaryhiddenRate.dobrushin_charpoly_rootMultiplicity_one

\item \emph{The coefficient modulus.}\label{proof:pair-polynomial-modulus:coefficient-modulus}
If \(q_P(1)=0\), the factorization of \(\chi_P\) would make it divisible by
\((z-1)^2\), contradicting \(m_P=1\).
The primed factor satisfies the same property. Consequently,
\[
q_P(1)\ne0,\qquad q_{P'}(1)\ne0,\qquad
|Q_{P,P'}(1)|=|q_P(1)|\,|q_{P'}(1)|>0.
\]
Define the seven moment differences and their maximum by
\[
\varepsilon_k:=\nu P^kr-\nu'P'^kr'
\quad(k\in\{0,\ldots,6\}),\qquad
\varepsilon_{\max}:=\max_{0\le k\le6}|\varepsilon_k|.
\]
Taking absolute values in the identity from \cref{proof:pair-polynomial-modulus:identity} gives
\[
\begin{aligned}
|Q_{P,P'}(1)|\,|\pi r-\pi'r'|
&=\left|\sum_{k=0}^6Q_k\varepsilon_k\right|\\
&\le\sum_{k=0}^6|Q_k|\,|\varepsilon_k|\\
&\le\left(\sum_{k=0}^6|Q_k|\right)\varepsilon_{\max}.
\end{aligned}
\]
Division by the positive denominator therefore yields
\[
|\pi r-\pi'r'|
\le
\frac{\sum_{k=0}^6|Q_k|}{|Q_{P,P'}(1)|}\,
\varepsilon_{\max}.
\]
% lean: CausalSmith.Stat.PomdpBinaryhiddenRate.pairPolynomial_value_coefficient_modulus

\item \emph{The complex root bound.}
Let \(z\in\mathbb C\) be a root of \(Q_{P,P'}\).
The product factorization gives the two cases
\[
q_P(z)=0\qquad\text{or}\qquad q_{P'}(z)=0.
\]
Consider the first case.
Since \(q_P(1)\ne0\), we have \(z\ne1\), while
\[
\chi_P(z)=(z-1)q_P(z)=0.
\]
Regard \(P\) as acting on complex column vectors.
Singularity of \(z\,\mathrm{Id}-P\) supplies
\[
f_{\mathbb C}\in\mathbb C^{\mathcal S}\setminus\{0\},
\qquad Pf_{\mathbb C}=zf_{\mathbb C}.
\]
A constant vector is fixed by \(P\), so a constant \(f_{\mathbb C}\) would satisfy
\[
(z-1)f_{\mathbb C}=0,
\]
forcing \(z=1\). Thus its diameter is positive:
\[
d_f:=\max_{i,j\in\mathcal S}
|f_{\mathbb C}(i)-f_{\mathbb C}(j)|>0.
\]

To bound this diameter after applying \(P\), take an arbitrary complex function and define
\[
\varphi:\mathcal S\to\mathbb C,\qquad
d_\varphi:=\max_{i,j\in\mathcal S}|\varphi(i)-\varphi(j)|,
\qquad
W_\varphi:=\sum_s(\omega_1(s)-\omega_2(s))\varphi(s),
\]
where \(\omega_1,\omega_2\) are probability vectors.
If \(W_\varphi=0\), then
\[
|W_\varphi|=0\le d_\varphi\tau(\omega_1,\omega_2),
\]
because both factors on the right are nonnegative.
If \(W_\varphi\ne0\), define the real function
\[
\gamma_\varphi:\mathcal S\to\mathbb R,\qquad
\gamma_\varphi(s):=
\operatorname{Re}\bigl(\overline{W_\varphi}\,\varphi(s)\bigr).
\]
For every \(i,j\in\mathcal S\),
\[
\begin{aligned}
|\gamma_\varphi(i)-\gamma_\varphi(j)|
&=\left|\operatorname{Re}
  \bigl(\overline{W_\varphi}(\varphi(i)-\varphi(j))\bigr)\right|\\
&\le |W_\varphi|\,|\varphi(i)-\varphi(j)|
\le |W_\varphi|\,d_\varphi.
\end{aligned}
\]
In particular,
\[
\max_s\gamma_\varphi(s)-\min_s\gamma_\varphi(s)
\le |W_\varphi|\,d_\varphi.
\]
The defining sum for \(W_\varphi\) also gives
\[
\sum_s(\omega_1(s)-\omega_2(s))\gamma_\varphi(s)
=\operatorname{Re}\bigl(\overline{W_\varphi}W_\varphi\bigr)
=|W_\varphi|^2.
\]
Applying the real expectation bound from \cref{proof:pair-polynomial-modulus:fixed-vectors} to
\(\gamma_\varphi\), therefore,
\[
\begin{aligned}
|W_\varphi|^2
&=\left|\sum_s(\omega_1(s)-\omega_2(s))\gamma_\varphi(s)\right|\\
&\le
\bigl(\max_s\gamma_\varphi(s)-\min_s\gamma_\varphi(s)\bigr)
\tau(\omega_1,\omega_2)\\
&\le |W_\varphi|\,d_\varphi\tau(\omega_1,\omega_2).
\end{aligned}
\]
Dividing by \(|W_\varphi|>0\), and including the zero case already treated, proves
\[
\left|\sum_s(\omega_1(s)-\omega_2(s))\varphi(s)\right|
\le d_\varphi\tau(\omega_1,\omega_2).
\]
% lean: Causalean.Mathlib.Probability.FiniteMarkovOscillation.probability_complex_mean_tv_diameter

Apply this estimate to two rows of \(P\). For every \(i,j\in\mathcal S\),
\[
\begin{aligned}
|(P\varphi)(i)-(P\varphi)(j)|
&=\left|\sum_s(P_{is}-P_{js})\varphi(s)\right|\\
&\le d_\varphi\tau\bigl(P(i,\cdot),P(j,\cdot)\bigr)\\
&\le d_\varphi\delta(P)
\le\alpha d_\varphi.
\end{aligned}
\]
Choose a pair attaining the diameter of \(f_{\mathbb C}\):
\[
i_{\mathrm{diam}},j_{\mathrm{diam}}\in\mathcal S,\qquad
|f_{\mathbb C}(i_{\mathrm{diam}})
 -f_{\mathbb C}(j_{\mathrm{diam}})|=d_f.
\]
Using the eigenvector equation in the preceding contraction estimate gives
\[
\begin{aligned}
|z|\,d_f
&=|zf_{\mathbb C}(i_{\mathrm{diam}})
     -zf_{\mathbb C}(j_{\mathrm{diam}})|\\
&=|(Pf_{\mathbb C})(i_{\mathrm{diam}})
     -(Pf_{\mathbb C})(j_{\mathrm{diam}})|\\
&\le\alpha d_f.
\end{aligned}
\]
Since \(d_f>0\), division gives \(|z|\le\alpha\).
In the second case, \(q_{P'}(z)=0\), the same argument uses \(P'\) and
\(\delta(P')\le\alpha\). We have therefore proved
\[
Q_{P,P'}(z)=0\quad\Longrightarrow\quad |z|\le\alpha
\qquad(z\in\mathbb C).
\]
% lean: CausalSmith.Stat.PomdpBinaryhiddenRate.pairPoly_complex_roots_norm_le

\item \emph{The coefficient ratio.}
Each characteristic polynomial is monic of degree four.
Its factorization by the monic linear polynomial \(z-1\) shows that
\[
q_P,\ q_{P'}\text{ are monic of degree }3,
\qquad
Q_{P,P'}\text{ is monic},\qquad
\deg Q_{P,P'}=3+3=6.
\]
Factoring over \(\mathbb C\), enumerate the roots with algebraic multiplicity:
\[
\xi_1,\ldots,\xi_6\in\mathbb C,\qquad
Q_{P,P'}(z)=\prod_{\ell=1}^6(z-\xi_\ell),
\qquad |\xi_\ell|\le\alpha\quad(1\le\ell\le6).
\]
For each \(k\in\{0,\ldots,6\}\), define
\[
\mathcal J_k
:=\{J\subseteq\{1,\ldots,6\}:\#J=6-k\}.
\]
Expanding the product, with empty products interpreted as one, gives
\[
Q_k=(-1)^{6-k}
      \sum_{J\in\mathcal J_k}\prod_{\ell\in J}\xi_\ell,
\qquad
\#\mathcal J_k=\binom6{6-k}=\binom6k.
\]
Thus every coefficient satisfies
\[
|Q_k|
\le\sum_{J\in\mathcal J_k}\prod_{\ell\in J}|\xi_\ell|
\le\binom6k\alpha^{6-k}.
\]
Summing these seven bounds and using the binomial formula yields
\[
\sum_{k=0}^6|Q_k|
\le\sum_{k=0}^6\binom6k\alpha^{6-k}
=(1+\alpha)^6.
\]
% lean: CausalSmith.Stat.PomdpBinaryhiddenRate.degreeSix_coeff_sum_le_of_roots_le

For the denominator, the reverse triangle inequality gives
\[
|1-\xi_\ell|
\ge1-|\xi_\ell|
\ge1-\alpha>0
\qquad(1\le\ell\le6).
\]
The empty product equals one. Successively multiplying the nonnegative
distance bounds gives, by induction,
\[
\prod_{\ell=1}^{n_{\mathrm{root}}}|1-\xi_\ell|
\ge(1-\alpha)^{n_{\mathrm{root}}}
\qquad
(n_{\mathrm{root}}\in\{0,\ldots,6\}).
\]
Indeed, for \(n_{\mathrm{root}}<6\), the induction step is
\[
\begin{aligned}
\prod_{\ell=1}^{n_{\mathrm{root}}+1}|1-\xi_\ell|
&=|1-\xi_{n_{\mathrm{root}}+1}|
       \prod_{\ell=1}^{n_{\mathrm{root}}}|1-\xi_\ell|\\
&\ge(1-\alpha)
       \prod_{\ell=1}^{n_{\mathrm{root}}}|1-\xi_\ell|\\
&\ge(1-\alpha)^{n_{\mathrm{root}}+1}.
\end{aligned}
\]
Evaluating the root factorization at one consequently gives
\[
|Q_{P,P'}(1)|
=\prod_{\ell=1}^6|1-\xi_\ell|
\ge(1-\alpha)^6>0.
\]
% lean: CausalSmith.Stat.PomdpBinaryhiddenRate.degreeSix_eval_one_lower_of_roots_le

Combining the numerator and denominator bounds, with positive denominators,
proves
\[
\frac{\sum_{k=0}^6|Q_k|}{|Q_{P,P'}(1)|}
\le\frac{(1+\alpha)^6}{(1-\alpha)^6}
=\left(\frac{1+\alpha}{1-\alpha}\right)^6
=B_\alpha.
\]
% lean: CausalSmith.Stat.PomdpBinaryhiddenRate.degreeSix_coefficient_ratio_le_of_roots_le

\item \emph{The value modulus.}
The maximum is nonnegative because its index set contains zero:
\[
0\le|\varepsilon_0|\le\varepsilon_{\max}.
\]
We may therefore multiply the coefficient-ratio bound by
\(\varepsilon_{\max}\). Together with \cref{proof:pair-polynomial-modulus:coefficient-modulus}, this gives
\[
\begin{aligned}
|\pi r-\pi'r'|
&\le
\frac{\sum_{k=0}^6|Q_k|}{|Q_{P,P'}(1)|}\,
\varepsilon_{\max}\\
&\le B_\alpha\varepsilon_{\max}\\
&=B_\alpha
  \max_{0\le k\le6}|\nu P^kr-\nu'P'^kr'|.
\end{aligned}
\]
% lean: CausalSmith.Stat.PomdpBinaryhiddenRate.pairPolynomial_value_modulus
\end{enumerate}
\end{proof}
